\documentclass[11pt,a4paper]{article}
\usepackage[margin=24mm]{geometry}
\usepackage[T1]{fontenc}
\usepackage[utf8]{inputenc}
\usepackage{lmodern}
\usepackage{microtype}
\usepackage{graphicx}
\usepackage{booktabs}
\usepackage{array}
\usepackage{amsmath,amssymb}
\usepackage{float}
\usepackage{caption}
\usepackage{hyperref}
\usepackage{xcolor}
\usepackage{enumitem}
\usepackage{fancyhdr}
\graphicspath{{artifacts/local/figures/}}
\hypersetup{colorlinks=true,linkcolor=blue!55!black,urlcolor=blue!55!black}
\setlength{\parskip}{6pt}
\setlength{\parindent}{0pt}
\setlist[itemize]{topsep=3pt,itemsep=2pt}
\setlength{\headheight}{14pt}
\pagestyle{fancy}
\fancyhf{}
\lhead{ID2223 Lab 1}
\rhead{CIFAR-10: Tasks 1--6}
\cfoot{\thepage}
\captionsetup{font=small,labelfont=bf}
\newcommand{\fig}[3]{\begin{figure}[htbp]\centering\includegraphics[width=#2\linewidth]{#1}\caption{#3}\end{figure}}

\begin{document}

\begin{titlepage}
\centering
\vspace*{28mm}
{\Large KTH Royal Institute of Technology\par}
\vspace{8mm}
{\large ID2223 Scalable Machine Learning and Deep Learning\par}
\vspace{20mm}
{\Huge\bfseries Lab 1: CIFAR-10 Classification\par}
\vspace{6mm}
{\LARGE AlexNet and ResNet-18: accuracy, robustness and cost\par}
\vspace{20mm}
{\large Report for Tasks 1--6\par}
\vfill
\begin{tabular}{rl}
Group: & Group Skibzz\\[3mm]
Students: & Md Ahsanul Karim, \texttt{makari@kth.se}\\
 & Md Sakibul Islam, \texttt{msis3@kth.se}\\
 & Rong Jin, \texttt{rongjin0715@gmail.com}\\[5mm]
Date: & \today\\[5mm]
Notebook: & \texttt{ID2223\_Lab1\_Full\_End\_to\_End.ipynb}\\
Hardware: & Google Colab, NVIDIA Tesla T4\\
\end{tabular}
\vspace{25mm}
\end{titlepage}

\setcounter{tocdepth}{1}
\tableofcontents
\newpage

% ---------------------------------------------------------------------------
\section{Introduction and setup}

We train two convolutional networks on CIFAR-10, a CIFAR-adapted AlexNet and a CIFAR-adapted ResNet-18, and compare them on accuracy, robustness and cost. Tasks 1 to 5 prepare the data, build and train the models, vary capacity and regularization, test robustness to corrupted inputs and look at individual predictions and adversarial examples. Task 6 puts all configurations side by side and asks which one we would actually pick.

Every number and figure in this report comes from a single top-to-bottom run of the notebook \texttt{ID2223\_Lab1\_\allowbreak Full\_\allowbreak End\_\allowbreak to\_\allowbreak End.ipynb} on a Google Colab Tesla T4. All models were retrained in that run (\texttt{FORCE\_RETRAIN = True}), so the timings are comparable with each other. The figures are the files the notebook saved to \texttt{artifacts/local/figures/}.

\paragraph{Compute budget.} We fixed one budget for every experiment: eight epochs over the 45,000 training images, batch size 64, SGD with learning rate 0.01 and momentum 0.9, no learning-rate schedule, and 32-bit floating point on one T4. Weight decay is $5\times10^{-4}$ and dropout 0.5 unless an experiment changes them. Under this budget one AlexNet run takes about 80 seconds and one ResNet-18 run about six minutes.

% ---------------------------------------------------------------------------
\section{Task 1: Study and prepare the data}

\subsection{What an example is, and how the classes are distributed}
A CIFAR-10 example is a $32\times32$ colour photograph stored as a $(3,32,32)$ tensor of 8-bit values in $[0,255]$, labelled with one of ten classes: airplane, automobile, bird, cat, deer, dog, frog, horse, ship and truck. The provided training set has 50,000 images and the test set 10,000. Each class has exactly 5,000 training and 1,000 test images, so overall accuracy is not skewed by class frequency. Per-class results still matter, as Task 5 shows.

\fig{t1_examples.png}{0.5}{Eight random training images per class. At $32\times32$ pixels, viewpoint, background and object size vary a lot, and some objects fill only a small part of the frame.}

Simple pixel statistics already show that some classes come with typical colours. Airplane and ship images are the brightest and bluest (mean blue intensity 0.589 and 0.555, against 0.446 for the whole training split), because sky and water fill much of the frame. Frog images are the darkest (mean brightness 0.418) and the least blue (0.345). A network can use such background cues, which helps on clean data and can hurt when the colours are disturbed.

\subsection{Split and preprocessing}
We split the 50,000 provided training images with a stratified random split (\texttt{SPLIT\_SEED = 2223}) into 45,000 for training and 5,000 for validation, which gives 4,500 and 500 images per class. The 10,000 test images stay untouched until final evaluation. Validation accuracy is used to pick checkpoints and compare settings; test accuracy is only reported.

Both networks get the same input: pixels are scaled to $[0,1]$ and then standardized per channel,
\[
x'_c = \frac{x_c-\mu_c}{\sigma_c},\qquad
\boldsymbol{\mu}=(0.4913,\,0.4822,\,0.4464),\quad
\boldsymbol{\sigma}=(0.2470,\,0.2434,\,0.2615).
\]
The mean and standard deviation are computed on the 45,000 training images only, so no information from validation or test leaks into preprocessing. No resizing is needed because both architectures are adapted to $32\times32$ inputs. The splits are kept on the GPU as 8-bit tensors and converted batch by batch, which keeps the data pipeline fast and reusable across all experiments. Each image keeps its original index, so predictions can be traced back to specific images.

Augmentation (used from Task 4 onwards) is applied in the $[0,1]$ pixel domain before normalization and only to training batches. \emph{Standard} augmentation pads by 4 pixels with reflection, crops back to $32\times32$ at a random position and flips horizontally with probability 0.5. \emph{Strong} augmentation adds, for about half of the images, a random rotation of up to $15^\circ$ and a shear of up to $10^\circ$, and for about half, Gaussian noise with a standard deviation drawn up to 0.06.

\fig{t1_augmentations.png}{0.7}{The same six training images without augmentation, with standard augmentation and with strong augmentation.}

\subsection{Controlling randomness}
Three seeds control the randomness: \texttt{SPLIT\_SEED = 2223} for the split, \texttt{BASE\_SEED = 42} for weight initialization, dropout and the shuffling order, and fixed seeds for the corrupted evaluation images (Task 4) and attacks (Task 5). Python, NumPy and PyTorch are seeded together, cuDNN benchmarking is off and deterministic algorithms are requested. Each training run also owns its own generator for shuffling and augmentation, so runs do not affect each other's random streams.

This works in practice. The Task 3 and Task 4 runs reached exactly the same validation accuracies, to two decimals, as an earlier run of the same code in a different Colab T4 session. Exact equality is not guaranteed on other GPUs or library versions, and timings always vary somewhat between sessions.

% ---------------------------------------------------------------------------
\section{Task 2: Implement and compare AlexNet and ResNet-18}

\subsection{AlexNet for CIFAR-10}
The original AlexNet was designed for $224\times224$ ImageNet images and starts with an $11\times11$ convolution with stride 4. On a $32\times32$ image that stem would throw away most of the spatial detail in the first layer, so our version uses five $3\times3$ convolutions with padding 1 and three $2\times2$ max-pooling steps (Table~\ref{tab:alexnet}). We keep the original pattern of five convolutional layers with channel widths 64, 192, 384, 256 and 256, followed by a fully connected classifier with dropout. The classifier is smaller than the original (1024 and 512 hidden units instead of 4096 and 4096), because the flattened feature vector has only 4,096 values. The network has 6{,}976{,}842 trainable parameters.

\begin{table}[htbp]\centering\small
\caption{CIFAR-10 AlexNet: output size of each stage for one image.}\label{tab:alexnet}
\begin{tabular}{ll}\toprule
Stage & Output\\\midrule
Input & $3\times32\times32$\\
Conv $3\times3$ (64), ReLU, max-pool & $64\times16\times16$\\
Conv $3\times3$ (192), ReLU, max-pool & $192\times8\times8$\\
Conv $3\times3$ (384), ReLU & $384\times8\times8$\\
Conv $3\times3$ (256), ReLU & $256\times8\times8$\\
Conv $3\times3$ (256), ReLU, max-pool & $256\times4\times4$\\
Flatten & 4096\\
Linear, ReLU, dropout 0.5 & 1024\\
Linear, ReLU, dropout 0.5 & 512\\
Linear & 10 logits\\\bottomrule
\end{tabular}\end{table}

Each convolution slides the same small set of learned filters over the whole image, so a pattern is detected wherever it appears. The first layer responds to edges, colour blobs and simple textures. Later layers combine these into parts and shapes, and every pooling step halves the resolution so that a $3\times3$ filter deeper in the network covers a larger part of the image. The fully connected head then maps the final $256\times4\times4$ feature map to ten class scores.

\subsection{ResNet-18 for CIFAR-10}
ResNet-18 replaces the original $7\times7$ stride-2 stem and the max-pool with a single $3\times3$ convolution with stride 1, again so that the small image is not downsampled too early. After the stem come four stages of two basic blocks each, with 64, 128, 256 and 512 channels at resolutions $32\times32$, $16\times16$, $8\times8$ and $4\times4$. Global average pooling reduces the last feature map to 512 numbers, and one linear layer produces the ten logits. The network has 11{,}173{,}962 trainable parameters, 1.60 times as many as our AlexNet.

A basic block computes two $3\times3$ convolutions with batch normalization and adds its input back before the final ReLU:
\[
\mathbf{y}=\operatorname{ReLU}\bigl(F(\mathbf{x})+S(\mathbf{x})\bigr).
\]
$S$ is the identity when the shape stays the same. When a stage halves the resolution and doubles the channels, $S$ is a $1\times1$ convolution with stride 2 followed by batch normalization. Because the identity path is always there, each block only has to learn a correction to its input, and gradients reach the early layers directly through the shortcuts. This is what makes networks of this depth trainable with plain SGD.

Both models passed a shape check: a batch of two images gives an output of shape $(2,10)$.

\subsection{Training and checkpoint selection}
Both baselines use the compute budget from the introduction with cross-entropy loss on the logits. After every epoch the model is evaluated on the 5,000 validation images with dropout off and without gradient updates. The general training routine saves the latest state every epoch (weights, optimizer and shuffling generator, so a run can resume) and keeps a separate copy of the epoch with the highest validation accuracy. That best checkpoint is the one used afterwards. The Task 2 AlexNet was trained with an explicit loop that validates every epoch and saves the final weights; its best validation epoch was also its last.

\subsection{Results}

\begin{table}[htbp]\centering\small
\caption{Task 2 baselines. Validation accuracy is from the selected checkpoint; test accuracy and loss are on the 10,000 official test images. AlexNet time is from the identical Task 3 baseline run, because the Task 2 loop did not record time.}\label{tab:baselines}
\begin{tabular}{lrrrrr}\toprule
Model & Parameters & Best val. acc. & Test acc. & Test loss & Training time\\\midrule
AlexNet & 6{,}976{,}842 & 78.04\% (epoch 8) & 77.95\% & 0.6767 & 78.9 s\\
ResNet-18 & 11{,}173{,}962 & 82.64\% (epoch 8) & 81.71\% & 0.6316 & 365.2 s\\\bottomrule
\end{tabular}\end{table}

\fig{t2_accuracy_comparison.png}{0.95}{Training and validation accuracy per epoch for the two baselines.}

ResNet-18 is 3.76 percentage points more accurate on the test set and also has the lower test loss. It learns much faster: after one epoch it is at 56.1\% validation accuracy, where AlexNet is at 41.5\%. It also fits the training data far more closely. At the last epoch the gap between training and validation accuracy is 12.7 points for ResNet-18 (95.3\% against 82.6\%) and 3.9 points for AlexNet (81.9\% against 78.0\%). So ResNet-18 overfits more, yet still generalizes better.

Its validation curve is also bumpier. It drops from 76.4\% at epoch 3 to 75.2\% at epoch 5 before climbing to 82.6\%. With a constant learning rate and batch normalization, single-epoch validation scores move around, so the choice of the ``best'' epoch is partly luck. Task 6 comes back to this.

The cost is real. One ResNet-18 epoch takes 45.7 seconds on the T4, against 9.9 seconds for AlexNet, so the full run takes 4.6 times as long.

% ---------------------------------------------------------------------------
\section{Task 3: Capacity and regularization}

All comparisons use the AlexNet baseline (classifier $1024\to512$, dropout 0.5, weight decay $5\times10^{-4}$) and change exactly one setting. The split, preprocessing, optimizer, learning rate, batch size, seed and eight-epoch budget stay fixed. This baseline is trained by the general training routine and differs from the Task 2 AlexNet only in its random state (initial weights and dropout masks). Its accuracy (78.56\% validation, 78.71\% test) is within the seed-to-seed variation measured in Task 6.

\subsection{Capacity: classifier width}
We keep the five convolutional layers and change only the width of the two hidden linear layers: small ($256\to128$), baseline ($1024\to512$) and large ($2048\to1024$).

\begin{table}[htbp]\centering\small
\caption{Capacity comparison after eight epochs.}
\begin{tabular}{lrrrr}\toprule
Width & Parameters & Final train acc. & Best val. acc. & Training time\\\midrule
Small & 3{,}334{,}602 & 77.20\% & 76.72\% & 72.4 s\\
Baseline & 6{,}976{,}842 & 82.22\% & 78.56\% & 78.9 s\\
Large & 12{,}750{,}666 & 83.42\% & 79.06\% & 88.8 s\\\bottomrule
\end{tabular}\end{table}

\fig{t3_capacity_accuracy.png}{0.95}{Training and validation accuracy for three classifier widths.}
\fig{t3_capacity_loss.png}{0.95}{Training and validation loss for three classifier widths.}

More width helps, but less and less. Going from small to baseline adds 1.84 points of validation accuracy; going from baseline to large adds only 0.50, for 1.83 times the parameters and 12.5\% more training time. The small model fits the training data least (77.2\%) and has the highest validation loss throughout, which points to some underfitting. None of the three shows clear overfitting yet. Validation loss is still falling for all of them at epoch 8, and all three reach their best validation accuracy in the last epoch. On the test set, the large model (78.68\%) is no better than the baseline (78.71\%).

\subsection{L2 regularization}
In PyTorch's SGD, weight decay adds $\lambda\mathbf{w}$ to every gradient, which is the same as an L2 penalty $\tfrac{\lambda}{2}\|\mathbf{w}\|^2$ on the loss. We compare $\lambda=0$, the baseline $5\times10^{-4}$ and a strong $5\times10^{-3}$.

\begin{table}[htbp]\centering\small
\caption{Weight decay comparison.}
\begin{tabular}{lrrrr}\toprule
Weight decay & Final train acc. & Best val. acc. & Best epoch & Training time\\\midrule
0 & 83.92\% & 79.94\% & 8 & 75.2 s\\
0.0005 & 82.22\% & 78.56\% & 8 & 78.9 s\\
0.005 & 67.64\% & 69.28\% & 8 & 79.0 s\\\bottomrule
\end{tabular}\end{table}

\fig{t3_l2_comparison.png}{0.85}{Training and validation accuracy and loss for three weight decay settings.}

Strong weight decay clearly hurts. Training accuracy only reaches 67.6\% and validation 69.3\%, so the penalty keeps the weights too small for the network to fit the data in eight epochs. That is underfitting, not better generalization. Removing weight decay gave the best validation accuracy (79.94\%, 1.38 points above the baseline). That difference is about the size of the seed-to-seed spread in Task 6, though, so we would not conclude that L2 regularization is useless. Over eight epochs the models are not yet overfitting much, and that is when weight decay would matter. Weight decay costs nothing extra to compute; the small time differences are run-to-run noise.

\subsection{Dropout}
We compare no dropout, dropout 0.3 and the baseline's 0.5, applied after both hidden linear layers.

\begin{table}[htbp]\centering\small
\caption{Dropout comparison. The train-validation gap is measured at the last epoch.}
\begin{tabular}{lrrrrr}\toprule
Dropout & Final train acc. & Final val. acc. & Gap & Best val. acc. (epoch) & Training time\\\midrule
0.0 & 85.37\% & 77.64\% & 7.73 & 78.64\% (7) & 78.9 s\\
0.3 & 83.53\% & 78.22\% & 5.31 & 78.46\% (7) & 79.4 s\\
0.5 & 82.22\% & 78.56\% & 3.66 & 78.56\% (8) & 78.9 s\\\bottomrule
\end{tabular}\end{table}

\fig{t3_dropout_comparison.png}{0.85}{Training and validation accuracy and loss for three dropout probabilities.}

Dropout does what it should. Without it the network fits the training data best (85.4\%), but validation accuracy peaks at epoch 7 and then falls by one point, the first sign of overfitting in this study. The train-validation gap shrinks from 7.7 points without dropout to 3.7 points with 0.5. The best validation accuracies are within 0.2 points of each other, so dropout changes how the model learns more than how good it ends up after eight epochs. Dropout has no measurable effect on training time.

\subsection{Cost of capacity and regularization}
Only width changes the cost. Parameters go from 3.3 to 7.0 to 12.8 million and training time from 72.4 to 78.9 to 88.8 seconds; the convolutions dominate the compute, so doubling the classifier adds much less than double the time. Weight decay and dropout leave the parameter count at 6{,}976{,}842 and the time unchanged. The logged peak GPU memory rises from 543 MB for the first run to over 900 MB for later ones, but that is because earlier models stay in memory during the session. These peaks are not per-model footprints.

% ---------------------------------------------------------------------------
\section{Task 4: Robustness}

\subsection{Training conditions and corrupted evaluation set}
We train the AlexNet baseline three ways: without augmentation (the Task 3 baseline), with standard augmentation (random crop and flip) and with strong augmentation (crop and flip plus random rotation, shear and noise; Section 2.2). Everything else is identical.

For evaluation we take the 5,000 validation images and build three fixed versions: the clean images, every image rotated by $15^\circ$ (bilinear, reflection padding), and every image with Gaussian noise of standard deviation 0.06 added on the $[0,1]$ scale (fixed seed, clipped to $[0,1]$). All models see exactly the same pixels, and no model is trained or tuned on these images. We picked these two corruptions because they mimic a tilted camera and sensor noise, and kept them moderate so the label stays obvious to a person.

\subsection{Clean and corrupted accuracy}

\begin{table}[htbp]\centering\small
\caption{Accuracy on the same 5,000 validation images. The last column is the drop from clean to noisy.}
\begin{tabular}{lrrrr}\toprule
Training condition & Clean & Rotation $15^\circ$ & Noise $\sigma=0.06$ & Noise drop\\\midrule
No augmentation & \textbf{78.56\%} & 70.32\% & 66.38\% & 12.18\\
Standard (crop, flip) & 77.56\% & 68.44\% & 65.06\% & 12.50\\
Strong (+ rotation, shear, noise) & 76.40\% & \textbf{71.86\%} & \textbf{73.74\%} & 2.66\\\bottomrule
\end{tabular}\end{table}

\fig{t4_robustness_comparison.png}{0.72}{Clean, rotated and noisy accuracy for the three training conditions.}

Strong augmentation trades 2.16 points of clean accuracy for 1.54 points under rotation and 7.36 points under noise. With it, adding noise costs the model only 2.7 points; without it, 12.2. Standard crop and flip does not help against these corruptions at all and is slightly worse everywhere after eight epochs. That makes sense: crops and flips teach invariance to position and mirroring, not to rotation or noise, and the extra variety slows down learning within a short budget. The gain from strong augmentation is largest for noise, which it trains on directly at up to the same strength. The evaluation rotation of $15^\circ$ sits at the edge of the training range, which may be why that gain is smaller.

\subsection{Which errors the corruptions introduce}
A new error is an image that a model classifies correctly when clean and wrongly after corruption. This is not the same as the drop in accuracy, because a few images also go from wrong to right.

\begin{table}[htbp]\centering\small
\caption{New errors among the 5,000 validation images.}
\begin{tabular}{lrr}\toprule
Corruption & No augmentation & Strong augmentation\\\midrule
Rotation $15^\circ$ & 647 (12.94\%) & 443 (8.86\%)\\
Gaussian noise & 732 (14.64\%) & 252 (5.04\%)\\\bottomrule
\end{tabular}\end{table}

\fig{t4_noise_error_examples.png}{0.88}{Images the unaugmented AlexNet gets right when clean and wrong after adding noise.}

Strong augmentation removes about two thirds of the noise-induced errors (732 down to 252) and a third of the rotation-induced ones (647 down to 443). The examples in the figure also show what kind of errors noise causes. Four of the five noisy images (a truck, a car, a deer and a dog) are classified as frog. Only one error stays within a related group (truck to automobile). Noise breaks up smooth surfaces into speckled texture, and frog images in CIFAR-10 are often dark, mottled close-ups, so the model seems to read the added texture as ``frog''. We did not plot rotation errors, so we cannot say whether rotation produces the same pattern.

\subsection{Cost}
Augmentation is free at inference time but not during training. Per epoch, training took 9.9 seconds without augmentation, 10.2 seconds with standard augmentation and 11.7 seconds with strong augmentation, so the full run costs 1.03 and 1.19 times the baseline. The extra time comes from the affine resampling and noise generation for every batch. Evaluating one model on 5,000 images takes about 0.34 seconds whichever way it was trained.

\subsection{Can strong transformations make labels ambiguous?}
Yes. At $32\times32$ pixels a few pixels decide what the object is. A large crop can cut off the part that defines the class (the wheels of a truck, the head of a horse), and a large rotation turns an airplane or ship into an image that looks like nothing in the training set. Very strong noise can hide a small animal entirely. Training on such images then teaches the model to predict a label that a person could no longer see. We therefore kept the training transformations moderate: rotations up to $15^\circ$, shears up to $10^\circ$ and noise up to 0.06, applied to only about half of the images. Our evaluation uses the same moderate strengths, so the results do not tell us how the models behave under stronger distortions.

% ---------------------------------------------------------------------------
\section{Task 5: Predictions and adversarial examples}

\subsection{Quantitative comparison on the test set}
We evaluate the two Task 2 baselines on the 10,000 test images.

\begin{table}[htbp]\centering\small
\caption{Test set metrics (macro averages over the ten classes).}
\begin{tabular}{lrrrr}\toprule
Model & Accuracy & Precision & Recall & F1\\\midrule
AlexNet & 77.9\% & 0.797 & 0.779 & 0.782\\
ResNet-18 & 81.7\% & 0.818 & 0.817 & 0.816\\\bottomrule
\end{tabular}\end{table}

\begin{table}[htbp]\centering\small
\caption{Per-class recall and F1 on the test set.}
\begin{tabular}{lrrrr}\toprule
& \multicolumn{2}{c}{AlexNet} & \multicolumn{2}{c}{ResNet-18}\\\cmidrule(lr){2-3}\cmidrule(lr){4-5}
Class & Recall & F1 & Recall & F1\\\midrule
airplane & 0.725 & 0.792 & 0.834 & 0.842\\
automobile & 0.914 & 0.890 & 0.922 & 0.909\\
bird & 0.674 & 0.710 & 0.676 & 0.732\\
cat & 0.726 & 0.608 & 0.653 & 0.666\\
deer & 0.863 & 0.757 & 0.779 & 0.779\\
dog & 0.624 & 0.689 & 0.755 & 0.734\\
frog & 0.865 & 0.836 & 0.911 & 0.844\\
horse & 0.761 & 0.824 & 0.877 & 0.871\\
ship & 0.872 & 0.880 & 0.913 & 0.898\\
truck & 0.771 & 0.838 & 0.851 & 0.887\\\bottomrule
\end{tabular}\end{table}

\fig{t5_confusion_matrices.png}{0.99}{Confusion matrices on the test set. Rows are true classes, columns predictions.}

Cat is the hardest class for both models (F1 0.608 and 0.666), followed by bird and dog. Vehicles are easy: automobile, ship and truck all have F1 around 0.84 to 0.91.

The confusion matrices show where the two models differ. AlexNet uses ``cat'' as a catch-all for animals: it predicts cat for 1,387 test images while only 1,000 are cats, which is why its cat precision is just 0.523. Its biggest single error is 244 dogs labelled cat, and it also sends 100 birds and 103 horses to deer. ResNet-18 spreads its errors more evenly. Its cat/dog confusion runs both ways (153 cats called dog, 127 dogs called cat), and its dog recall is 13 points higher than AlexNet's. Among vehicles, truck to automobile is the main error for both (97 and 64 images).

\subsection{Correct and incorrect predictions}

\fig{t5_prediction_categories.png}{0.85}{Test images grouped by outcome. Rows from top: both models correct, both wrong, only AlexNet correct, only ResNet-18 correct. Each title gives the true class and each model's prediction with its softmax confidence.}

We show the four combinations of correct and wrong, so that neither model is shown only at its worst. When both models fail, the image is often hard for a person too: a dog standing in grass that both models call deer, a close-up of a bird's head, a black-and-white photo of an old car that both call truck. Several mistakes come with high confidence, for example 100\% for deer on the dog image and 92\% for ship on an airplane that ResNet-18 gets right. Softmax confidence is therefore not a reliable sign that a prediction is correct.

\subsection{Adversarial examples with FGSM}
We use the fast gradient sign method, a one-step white-box attack that knows the model and its weights:
\[
\mathbf{x}_{\mathrm{adv}}=\operatorname{clip}_{[0,1]}\Bigl(\mathbf{x}+\varepsilon\,\operatorname{sign}\bigl(\nabla_{\mathbf{x}}\,\mathcal{L}(f_\theta(\mathbf{x}),y)\bigr)\Bigr).
\]
The gradient is taken with respect to the input pixels in $[0,1]$, through the normalization layer, while the weights stay frozen. Every pixel moves by at most $\varepsilon$, so the constraint is $\|\mathbf{x}_{\mathrm{adv}}-\mathbf{x}\|_\infty\le\varepsilon$. A check on eight images confirmed a maximum change of 0.0300 at $\varepsilon=0.03$. Each model is attacked with its own gradients on the same first 500 test images.

\begin{table}[htbp]\centering\small
\caption{Accuracy on 500 test images under FGSM.}
\begin{tabular}{lrrrr}\toprule
Model & $\varepsilon=0$ & $0.01$ & $0.03$ & $0.05$\\\midrule
AlexNet & 78.8\% & 34.6\% & 4.6\% & 0.8\%\\
ResNet-18 & 82.4\% & 16.0\% & 3.2\% & 1.6\%\\\bottomrule
\end{tabular}\end{table}

\fig{t5_fgsm_examples.png}{0.7}{Four test images that AlexNet classifies correctly, attacked with increasing $\varepsilon$. Titles give the prediction and its confidence.}

Both models break down almost completely. At $\varepsilon=0.03$, a change of under 8 intensity levels out of 255 per pixel, accuracy falls below 5\% for both. ResNet-18 starts higher but falls faster: at $\varepsilon=0.01$ it keeps only 16.0\% against 34.6\% for AlexNet. The attacked images in the figure look almost unchanged, but the model becomes more confident in the wrong answer as $\varepsilon$ grows. The cat becomes a dog with 94.3\% confidence at $\varepsilon=0.03$, and the first ship becomes an automobile with 99.8\%. At $\varepsilon=0.01$ the attack does not always succeed: two of the four examples (the cat and the frog) keep their correct label, and AlexNet still classifies 34.6\% of the 500 images correctly. From $\varepsilon=0.03$ on, all four examples are misclassified.

\subsection{Random corruption versus an optimized attack}
Gaussian noise in Task 4 is drawn without looking at the model. With $\sigma=0.06$ the typical pixel change is larger than FGSM's 0.03, yet the unaugmented AlexNet baseline still classifies 66.4\% of noisy images correctly. FGSM at half that per-pixel size leaves the same model at 6.8\% (Table~\ref{tab:tradeoff}). The difference is direction: FGSM moves every pixel the way that increases the loss most. Robustness to random corruption and robustness to an adversary are therefore different properties, and Task 6 shows that one does not predict the other.

% ---------------------------------------------------------------------------
\section{Task 6: Benchmark and tradeoffs}

Task 6 compares all ten configurations from Tasks 2 to 4 under the compute budget from the introduction: the AlexNet baseline and its eight variants, and ResNet-18. The AlexNet baseline here is the Task 3 run, because it was timed by the training routine. Every configuration is evaluated with the same functions on the same images: the test set, the Task 4 clean, rotated and noisy validation images, and FGSM at $\varepsilon=0.03$ on the 500 Task 5 images.

\subsection{Training cost and inference performance}
Training time is the sum of the logged epoch times, which include the validation pass. Inference is measured separately on the T4 in FP32: latency is the mean time for one image (batch size 1, 200 repeats), throughput uses batches of 256 (50 repeats). Both exclude five warm-up passes and synchronize the GPU before reading the clock. Weight decay, dropout and augmentation do not change the network, so only the four distinct architectures need to be timed.

\begin{table}[htbp]\centering\small
\caption{Inference cost per architecture on a Tesla T4 (FP32).}\label{tab:inference}
\begin{tabular}{lrrrr}\toprule
Architecture & Parameters & Weights & Latency (batch 1) & Throughput (batch 256)\\\midrule
AlexNet small & 3{,}334{,}602 & 13.3 MB & 1.04 ms & 15{,}624 img/s\\
AlexNet & 6{,}976{,}842 & 27.9 MB & 0.65 ms & 14{,}657 img/s\\
AlexNet large & 12{,}750{,}666 & 51.0 MB & 0.63 ms & 13{,}494 img/s\\
ResNet-18 & 11{,}173{,}962 & 44.8 MB & 2.61 ms & 2{,}149 img/s\\\bottomrule
\end{tabular}\end{table}

ResNet-18 is 4.0 times slower per image and has 6.8 times lower throughput than AlexNet. Its 11 million parameters are fewer than the large AlexNet's 12.8 million, yet it is far slower: its convolutions run on full $32\times32$ feature maps with up to 512 channels, while most of AlexNet's parameters sit in cheap fully connected layers. Parameter count is a poor guide to speed. The latency of the small AlexNet (1.04 ms) is higher than that of the larger ones, which cannot come from its compute. At batch size 1 the GPU is limited by kernel launch overhead, and the small model was timed first, so we treat differences below about half a millisecond as measurement noise.

\begin{table}[htbp]\centering\small
\caption{All configurations under the same budget. $\Delta$ columns are relative to the AlexNet baseline; time ratio is training time relative to the baseline.}\label{tab:tradeoff}
\setlength{\tabcolsep}{4pt}
\begin{tabular}{lrrrrrrrr}\toprule
Configuration & Train time & Ratio & Test & Rotation & Noise & FGSM & $\Delta$ test & $\Delta$ noise\\\midrule
AlexNet small & 72.4 s & 0.92 & 76.69 & 64.48 & 67.10 & 4.8 & $-2.02$ & $+0.72$\\
AlexNet baseline & 78.9 s & 1.00 & 78.71 & 70.32 & 66.38 & 6.8 & 0.00 & 0.00\\
AlexNet large & 88.8 s & 1.13 & 78.68 & 68.64 & 67.30 & 5.6 & $-0.03$ & $+0.92$\\
AlexNet no L2 & 75.2 s & 0.95 & 79.67 & 69.82 & 70.68 & 6.0 & $+0.96$ & $+4.30$\\
AlexNet strong L2 & 79.0 s & 1.00 & 68.69 & 58.94 & 63.30 & 6.2 & $-10.02$ & $-3.08$\\
AlexNet no dropout & 78.9 s & 1.00 & 76.97 & 66.64 & 70.08 & 5.2 & $-1.74$ & $+3.70$\\
AlexNet dropout 0.3 & 79.4 s & 1.01 & 77.02 & 65.66 & 69.90 & 6.4 & $-1.69$ & $+3.52$\\
AlexNet standard aug & 81.2 s & 1.03 & 77.45 & 68.44 & 65.06 & 5.6 & $-1.26$ & $-1.32$\\
AlexNet strong aug & 93.9 s & 1.19 & 76.18 & 71.86 & \textbf{73.74} & \textbf{9.6} & $-2.53$ & $+7.36$\\
ResNet-18 & 365.2 s & 4.63 & \textbf{81.71} & \textbf{71.86} & 62.98 & 3.2 & $+3.00$ & $-3.40$\\\bottomrule
\end{tabular}
\par\smallskip\footnotesize All accuracies in \%. Test on 10,000 images; rotation and noise on the 5,000 Task 4 validation images; FGSM ($\varepsilon=0.03$) on 500 test images.
\end{table}

\fig{t6_tradeoffs.png}{0.99}{Each configuration's test accuracy against training time (left) and inference throughput (middle), and its accuracy on noisy images against training time (right). AlexNet variants that share an architecture have the same throughput.}

Table~\ref{tab:tradeoff} gives the main picture. ResNet-18 is the most accurate on clean test images and, tied with strong augmentation, under rotation. But it is the least robust to noise of all ten configurations: its accuracy falls by 19.7 points when noise is added, against 12.2 for the AlexNet baseline and 2.7 for AlexNet with strong augmentation. It is also the easiest to fool with FGSM. The strongly augmented AlexNet is the most robust to both noise and FGSM, at the lowest clean accuracy among the reasonable configurations.

\subsection{Sensitivity to training duration, seed and hardware}

\paragraph{Training duration.} Eight of the ten configurations reached their best validation accuracy in the last epoch, so most of them were still improving when training stopped. Ranking the configurations by validation accuracy at epoch 4 and at epoch 8 gives the same top three (ResNet-18, no L2, large) and the same bottom two (strong augmentation, strong L2). Four of the five configurations in between change rank (the baseline moves from sixth to fourth, no dropout from fourth to sixth), so their order depends on when we stop.

\paragraph{Random seed.} We trained the AlexNet baseline and ResNet-18 twice more, with seeds 43 and 44. The seed changes the initial weights, dropout masks and shuffling order; the data split stays fixed.

\begin{table}[htbp]\centering\small
\caption{Three seeds per model (mean $\pm$ standard deviation, with range).}\label{tab:seeds}
\begin{tabular}{lccc}\toprule
Model & Best val. acc. & Test acc. & Training time\\\midrule
AlexNet baseline & $78.19\pm0.32$ (77.96--78.56) & $77.58\pm1.21$ (76.30--78.71) & $81.8\pm2.6$ s\\
ResNet-18 & $81.44\pm1.40$ (79.90--82.64) & $80.71\pm1.15$ (79.46--81.71) & $379.7\pm12.5$ s\\\bottomrule
\end{tabular}\end{table}

\fig{t6_seed_spread.png}{0.5}{Test accuracy for seeds 42, 43 and 44. The horizontal line is the mean.}

Test accuracy varies by about 1.2 points (standard deviation) between seeds for both models. ResNet-18 beats AlexNet for every seed, by 3.00, 3.24 and 3.16 points, so that difference is solid. Differences of about one point between single runs, which covers most of Task 3, are not. ResNet-18's own training was not stable either: with seed 43 its validation accuracy fell from 81.8\% at epoch 7 to 77.4\% at epoch 8, and its best epoch was 7 for two of the three seeds.

\paragraph{Hardware.} All timings come from one Tesla T4 in FP32. The notebook defines a mixed-precision flag, but the training and evaluation loops never use autocast, so no run used FP16. Even on the same GPU type, the seed runs of ResNet-18 took about 48 seconds per epoch against 45.7 for the first run, a 6\% difference with identical code. On another GPU, with mixed precision, or on a CPU or Apple MPS device, both the absolute times and the AlexNet-to-ResNet ratio would change. Accuracy depends much less on hardware: the Task 3 and 4 results were identical to two decimals across two T4 sessions.

\subsection{Discussion}

\paragraph{Which model would we choose, and why?}
Under our budget we would choose ResNet-18 when clean accuracy is what matters. It is 3.0 points more accurate than AlexNet on the test set, and it wins for all three seeds by more than the seed spread. A full training run takes about six minutes on a free Colab T4, which is affordable, and 2.6 ms per image is fast enough for most offline uses.

The choice changes if the inputs are noisy or inference has to be cheap. ResNet-18 loses 19.7 points under Gaussian noise, the most of any configuration, and runs at one seventh of AlexNet's throughput. AlexNet trained with strong augmentation costs 1.19 times the baseline's training time, nothing extra at inference, and keeps 73.7\% accuracy on noisy images against 63.0\% for ResNet-18. In that setting it is the better choice, at the price of about 5.5 points of clean test accuracy. Neither model resists FGSM (below 10\% at $\varepsilon=0.03$ for every configuration), so adversarial robustness does not decide between them.

\paragraph{Which changes improve accuracy or robustness at an acceptable cost?}
\begin{itemize}
\item Strong augmentation is the only change with a large, clear benefit: $+7.36$ points under noise and $+1.54$ under rotation for 19\% more training time and no inference cost. It costs 2.16 points of clean validation accuracy (2.53 on test). Worth it whenever corrupted inputs are expected.
\item Switching to ResNet-18 buys 3.0 points of clean accuracy for 4.6 times the training time and 6.8 times lower throughput, and makes the model less robust to noise. Worth it only if clean accuracy is the goal.
\item A wider classifier is not worth it: 1.83 times the parameters and 13\% more training time for no gain on the test set ($-0.03$ points).
\item Removing weight decay ($+0.96$ test, $+4.30$ noise) and removing or reducing dropout ($+3.5$ to $+3.7$ noise, about $-1.7$ test) are free, but the clean accuracy changes are within the seed spread, and the noise gains come from single runs. We would not adopt them without repeating them over seeds.
\item Strong weight decay costs nothing to compute but loses 10 points of accuracy. Standard crop and flip makes every metric slightly worse at eight epochs.
\end{itemize}

\paragraph{Which conclusions are sensitive to training duration, seed or hardware?}
The conclusions that hold up are ResNet-18's clean-accuracy lead (all three seeds), the robustness gain from strong augmentation (7.4 points, several times the seed spread), and the damage from strong weight decay (10 points). Everything within about 1.5 points is sensitive to the seed: the width comparison, the weight decay and dropout rankings, and the small noise differences between them. Most results also depend on training duration. Eight of ten configurations were still improving at epoch 8, and longer training usually helps ResNet-18 with batch normalization and augmented models more than the baseline, so the gaps could grow or shrink. All timings and time ratios belong to a T4 in FP32 and vary by about 6\% even between runs on the same GPU type.

\paragraph{What further experiment would give the most useful evidence?}
Train ResNet-18 with strong augmentation, and repeat it together with the AlexNet baseline, AlexNet with strong augmentation and plain ResNet-18 over at least three seeds with a longer schedule, for example 30 epochs with learning-rate decay. This one experiment would show whether strong augmentation fixes ResNet-18's weakness to noise, and whether the accuracy and robustness gaps survive longer training. That is exactly what our recommendation depends on. Evaluating the same models with a multi-step attack such as PGD would then tell us whether the small FGSM differences mean anything.

% ---------------------------------------------------------------------------
\section{Conclusions}
On CIFAR-10 with an eight-epoch budget on one T4, ResNet-18 reaches 81.7\% test accuracy against 78.0\% for our CIFAR-adapted AlexNet, and the lead holds across seeds. It costs 4.6 times the training time and about 4 times the per-image latency. Wider classifiers and changes to weight decay or dropout move accuracy by about a point, which is within the variation between seeds. Strong augmentation with rotation, shear and noise is the one cheap change that clearly improves robustness, cutting noise-induced errors from 732 to 252 on 5,000 images. None of the models resists a simple gradient attack, and better accuracy on clean images does not mean better robustness: the most accurate model was the most fragile under noise.

\section*{Reproduction}
\addcontentsline{toc}{section}{Reproduction}
Open \texttt{ID2223\_Lab1\_\allowbreak Full\_\allowbreak End\_\allowbreak to\_\allowbreak End.ipynb} in Google Colab, select a T4 GPU runtime and run all cells. The notebook downloads CIFAR-10 with torchvision, mounts Google Drive and stores checkpoints and figures under \texttt{MyDrive/kth-sml-lab1/local/}. With \texttt{FORCE\_RETRAIN = True} every model is retrained (about one hour on a T4); with \texttt{False}, finished runs are loaded from Drive and interrupted runs resume. Part A of the notebook contains the introductory Lab 1 tutorial and is independent of the CIFAR-10 experiments in Part B.

\section*{References}
\addcontentsline{toc}{section}{References}
\begin{enumerate}[leftmargin=*]
\item A. Krizhevsky. \emph{Learning Multiple Layers of Features from Tiny Images}. Technical report, University of Toronto, 2009. (CIFAR-10)
\item A. Krizhevsky, I. Sutskever and G. E. Hinton. ImageNet classification with deep convolutional neural networks. In \emph{NeurIPS}, 2012. (AlexNet)
\item K. He, X. Zhang, S. Ren and J. Sun. Deep residual learning for image recognition. In \emph{CVPR}, 2016. (ResNet)
\item I. J. Goodfellow, J. Shlens and C. Szegedy. Explaining and harnessing adversarial examples. In \emph{ICLR}, 2015. (FGSM)
\item S. Ioffe and C. Szegedy. Batch normalization: accelerating deep network training by reducing internal covariate shift. In \emph{ICML}, 2015.
\item N. Srivastava, G. Hinton, A. Krizhevsky, I. Sutskever and R. Salakhutdinov. Dropout: a simple way to prevent neural networks from overfitting. \emph{JMLR} 15, 2014.
\item PyTorch and torchvision documentation, \url{https://pytorch.org/docs/}. Course notebook \texttt{ID2223\_Lab1.ipynb}.
\end{enumerate}

\end{document}
